\documentclass[11pt,a4paper]{ctexart}
\usepackage[a4paper,margin=2cm]{geometry}
\usepackage{amsmath,amssymb}
\usepackage{booktabs}
\usepackage{longtable}
\usepackage{array}
\usepackage{enumitem}
\usepackage{setspace}
\usepackage{titlesec}
\usepackage{hyperref}
\usepackage{fontspec}
\setmainfont{Times New Roman}

\setstretch{1.2}

\titleformat{\section}{\centering\bfseries\large}{}{0pt}{}
\titleformat{\subsection}{\bfseries\normalsize}{}{0pt}{}

\begin{document}

\begin{center}
    {\LARGE \textbf{Introduction to Mao Zedong Thought and the Theoretical System of Socialism with Chinese Characteristics}}
\end{center}

\section{1919: The New Democratic Revolution}

\begin{enumerate}[leftmargin=2em]
    \item \textbf{The ``Two Integrations''}: Integrating the basic principles of Marxism (the ideological soul) with China's concrete realities, and with China's fine traditional culture (the cultural root).
    
    \item \textbf{Two stages of integration.}
    \begin{enumerate}[label=(\arabic*),leftmargin=2em]
        \item In 1938, Marxism was integrated with the Chinese revolution. In \textit{On the New Stage}, delivered at the Sixth Plenary Session of the Sixth Central Committee of the Communist Party of China, Mao Zedong put forward the proposition of the Sinicization of Marxism.
        \item In 1956, Marxism was integrated with socialist construction. This exploration was not completed.
    \end{enumerate}
    
    \item \textbf{The three magic weapons for defeating the enemy in the Chinese revolution}: The united front, armed struggle, and Party building.
    
    \item \textbf{The living soul of Mao Zedong Thought}: Seeking truth from facts, the mass line, and independence and self-reliance.
    
    \item \textbf{The connotations of seeking truth from facts}: Proceeding from reality in all respects, integrating theory with practice, and testing and developing truth through practice.
    
    \item \textbf{The connotations of the mass line}: Everything for the masses, everything relying on the masses, from the masses to the masses.
    
    \item \textbf{The connotations of independence and self-reliance}: Maintaining independent thinking and following one's own path.
    
    \item \textbf{The principal contradictions in modern China}: The contradiction between imperialism and the Chinese nation (the most important one), and the contradiction between feudalism and the masses of the people.
    
    \item \textbf{The three fundamental issues of the Chinese revolution}: The leadership of the proletariat (the central issue), the peasant question (the basic issue), and distinguishing between enemies and friends (the primary issue).
    
    \item \textbf{The motive forces of the New Democratic Revolution}: The working class (the most fundamental force), the peasantry (the main force), the urban petty bourgeoisie (allies), and the national bourgeoisie.
    
    \item \textbf{The general line of the New Democratic Revolution}: A revolution led by the proletariat, participated in by the masses, and directed against imperialism, feudalism, and bureaucrat-capitalism.
    
    \item \textbf{The objects of the New Democratic Revolution}: Imperialism (the primary target), feudalism, and \\ bureaucrat-capitalism.
    
    \item \textbf{The nature of the New Democratic Revolution}: A bourgeois-democratic revolution.
    
    \item \textbf{The politics of the New Democratic Revolution.}
    \begin{enumerate}[label=(\arabic*),leftmargin=2em]
        \item State system: Joint dictatorship of all revolutionary classes.
        \item Political system: The system of people's congresses under democratic centralism.
    \end{enumerate}
    
    \item \textbf{The general program of the New Democratic Revolution.}
    \begin{enumerate}[label=(\arabic*),leftmargin=2em]
        \item \textbf{Political program}: Overthrow the rule of imperialism and feudalism, and establish a New Democratic republic under the leadership of the proletariat, based on the worker-peasant alliance, and under the joint dictatorship of all revolutionary classes.
        \item \textbf{Economic program}: Confiscate the land of the feudal landlord class and distribute it to the peasants (the main content); confiscate the monopoly capital of the bureaucrat-bourgeoisie and turn it over to the New Democratic state; protect national industry and commerce.
        \item \textbf{Cultural program}: New Democratic culture is an anti-imperialist and anti-feudal culture of the masses under the leadership of the proletariat, namely a national, scientific, and mass culture.
    \end{enumerate}
    
    \item \textbf{Difference between the New Democratic Revolution and the socialist revolution}: They differ in nature. The New Democratic Revolution still belongs to the category of bourgeois-democratic revolution, whereas the socialist revolution is proletarian in nature; it aims to eliminate the capitalist system of exploitation and transform the private ownership of small-scale production, and politically to establish the dictatorship of the proletariat.
    
    \item \textbf{Connection between the New Democratic Revolution and the socialist revolution}: They are interconnected and closely linked. The New Democratic Revolution is the necessary preparation for the socialist revolution, and the socialist revolution is the inevitable trend of the New Democratic Revolution.
    
    \item \textbf{The Chinese revolutionary road of encircling the cities from the countryside and seizing political power by armed force}: Its essence lies in correctly handling the relationship among land revolution, armed struggle, and the building of rural revolutionary base areas.
\end{enumerate}

\section{II. 1949~Socialist Revolution and Initial Socialist Construction}

\begin{enumerate}[leftmargin=2em]
    \item \textbf{The Party's general line and general task during the transitional period}: ``One transformation'' and ``Three transformations,'' namely socialist industrialization, and the socialist transformation of agriculture, handicrafts, and capitalist industry and commerce.
    
    \item \textbf{The five economic sectors in the transitional period}: The state-owned economy, the cooperative economy, the individual economy, the private capitalist economy, and the state capitalist economy.
    
    \item \textbf{Significance of the establishment of the basic socialist system}: It marked the most profound and greatest social transformation in Chinese history, laid the institutional foundation for all development and progress in contemporary China, and provided an important prerequisite for the innovation and development of the socialist system with Chinese characteristics.
    
    \item \textbf{Independently exploring a socialist construction path suited to China's conditions}: This was put forward by Mao Zedong in \textit{On the Ten Major Relationships} at an enlarged meeting of the Political Bureau of the CPC Central Committee.
    
    \item \textbf{The Eighth National Congress of the Communist Party of China in 1956}: The principal contradiction was identified as the contradiction between the people's demand for building an advanced industrial country and the reality of a backward agricultural country, and the contradiction between the people's need for rapid economic and cultural development and the current inability of the economy and culture to satisfy those needs.
    
    \item \textbf{The Sixth Plenary Session of the Eleventh Central Committee in 1981}: The principal contradiction was defined as the contradiction between the ever-growing material and cultural needs of the people and backward social production.
    
    \item \textbf{The 19th National Congress of the Communist Party of China in 2017}: The principal contradiction was defined as the contradiction between the people's ever-growing needs for a better life and unbalanced and inadequate development.
    
    \item \textbf{``Two participations, one reform, and three-in-one combination''}: Cadres participating in labor, workers participating in management, reforming irrational rules and regulations, and combining workers, leading cadres, and technical personnel.
\end{enumerate}

\section{III. 1978~Deng Xiaoping Theory}

\begin{enumerate}[leftmargin=2em]
    \item \textbf{The central question of the times}: What is socialism and how to build socialism.
    
    \item \textbf{The essence of Deng Xiaoping Theory}: Emancipating the mind and seeking truth from facts.
    
    \item \textbf{The theme report of the Third Plenary Session of the Eleventh Central Committee}: \textit{Emancipate the Mind, Seek Truth from Facts, and Unite as One in Looking to the Future}.
    
    \item \textbf{The Party's basic line in the primary stage of socialism}: To lead and unite the people of all ethnic groups throughout the country, take economic development as the central task, uphold the Four Cardinal Principles, adhere to reform and opening up, rely on self-reliance and hard work, and strive to build China into a prosperous, strong, democratic, and culturally advanced socialist modernized country.
    
    \item \textbf{The essence of socialism}: To emancipate and develop the productive forces, eliminate exploitation, eliminate polarization, and ultimately achieve common prosperity.
    
    \item \textbf{How to build socialism}: One central task and two basic points, namely taking economic development as the central task, upholding the Four Cardinal Principles, and adhering to reform and opening up.
    
    \item \textbf{The Four Cardinal Principles}: Upholding the socialist road, the dictatorship of the proletariat, the leadership of the Communist Party, and Marxism-Leninism and Mao Zedong Thought.
    
    \item \textbf{The Four Modernizations}: Modernization of industry, agriculture, national defense, and science and technology.
    
    \item \textbf{The fundamental principle for socialist modernization}: Grasp both material progress and ideological-cultural progress, and both must be pursued with equal firmness.
\end{enumerate}

\section{IV. 1989~The Important Thought of Three Represents}

\begin{enumerate}[leftmargin=2em]
    \item \textbf{The central questions of the times}: Further answering what is socialism and how to build socialism, and also what kind of Party to build and how to build it.
    
    \item \textbf{The essence of the Important Thought of Three Represents}: Emancipating the mind, seeking truth from facts, and advancing with the times.
    
    \item \textbf{Main contents of the Important Thought of Three Represents}: The Communist Party of China must always represent.
    \begin{enumerate}[label=(\arabic*),leftmargin=2em]
        \item The development requirements of China's advanced productive forces.
        \item The forward direction of China's advanced culture.
        \item The fundamental interests of the overwhelming majority of the Chinese people.
    \end{enumerate}
    
    \item \textbf{To implement the Important Thought of Three Represents}: The key lies in advancing with the times, the core lies in maintaining the Party's progressiveness, and the essence lies in governing for the people.
    
    \item \textbf{The Important Thought of Three Represents}: The foundation of the Party, the cornerstone of governance, and the source of strength.
    
    \item \textbf{To establish a socialist market economy}: It is necessary to uphold and improve the basic socialist economic system in which public ownership remains the mainstay while diverse forms of ownership develop together.
    
    \item It is necessary to unswervingly consolidate and develop the public sector of the economy, and unswervingly encourage, support, and guide the development of the non-public sector.
\end{enumerate}

\section{V. 2002~The Scientific Outlook on Development}

\begin{enumerate}[leftmargin=2em]
    \item \textbf{The central question of the times}: Under new circumstances, what kind of development to achieve and how to achieve it.
    
    \item \textbf{The essence of the Scientific Outlook on Development}: Emancipating the mind, seeking truth from facts, advancing with the times, and being realistic and pragmatic.
    
    \item \textbf{The Scientific Outlook on Development}: Development is the top priority, putting people first is the core, comprehensive, balanced, and sustainable development is the basic requirement, and overall planning with all factors taken into consideration is the fundamental approach.
    
    \item \textbf{The three basic goals of all-round socialist modernization}: To build socialist political civilization, material civilization, and spiritual civilization.
    
    \item \textbf{Socialist democratic politics}: The leadership of the Party as the fundamental guarantee, the people being the masters of the country as the essential feature, and governing the country according to law as the basic strategy, all unified in the practice of socialist democratic politics.
\end{enumerate}

\section{Political Terms}

\begin{center}
\small
\setlength{\tabcolsep}{10pt}
\begin{tabular}{llll}
\toprule
Word & Part of Speech & IPA & Chinese \\
\midrule
proletariat & n. & /ˌprəʊlɪˈteəriət/ & 无产阶级 \\
working class & n. & /ˈwɜːkɪŋ klɑːs/ & 工人阶级 \\
peasantry & n. & /ˈpezəntri/ & 农民阶级 \\
urban petty bourgeoisie & n. & /ˈɜːbən ˌpeti ˌbʊəʒwɑːˈziː/ & 城市小资产阶级 \\
national bourgeoisie & n. & /ˈnæʃənəl ˌbʊəʒwɑːˈziː/ & 民族资产阶级 \\
feudal landlord class & n. & /ˈfjuːdəl ˈlændlɔːd klɑːs/ & 封建地主阶级 \\
bureaucrat-bourgeoisie & n. & /bjʊəˈrɒkræt ˌbʊəʒwɑːˈziː/ & 官僚资产阶级 \\
revolutionary classes & n. & /ˌrevəˈluːʃənəri ˈklɑːsɪz/ & 革命阶级 \\
\bottomrule
\end{tabular}
\end{center}

\end{document}